\documentclass[12pt]{article}
\usepackage[utf8]{inputenc}
\usepackage[T1]{fontenc}
\usepackage{hyperref}
\title{Koromu Language Resources for K-12 Curriculum Development}
\author{Compiled by Claude Sonnet 5.5 (claude-sonnet-5-5) for Andr\'as Kornai}
\date{}
\begin{document}
\maketitle

\section{Introduction}
Koromu, also called Kesawai (ISO 639-3 xes; Glottocode kesa1244), is a Rai Coast language of Madang Province, Papua New Guinea, in the Ramu Valley. Wikipedia cites about 700 native speakers (2008). Joshua Project reports about 1,500. Joshua Project also says younger people increasingly favour Tok Pisin. Only academic descriptions were found; no school materials turned up.

\section{Primary Educational Resources}
None were found.

\section{Linguistic and Reference Materials}
\begin{itemize}
\item Carol Priestley, \emph{Koromu (Kesawai): Grammar and Information Structure of a New Guinea Language} (Pacific Linguistics, 2019/2020). ANU research record: \url{https://researchportalplus.anu.edu.au/en/publications/koromu-kesawai-grammar-and-information-structure-of-a-new-guinea-/}
\item Glottolog reference for Priestley 2020: \url{https://glottolog.org/resource/reference/id/581512}
\item Glottolog reference for Priestley 2008 (the ANU PhD thesis on which the grammar is based): \url{https://glottolog.org/resource/reference/id/309219}
\item Wikipedia, Kesawai language: \url{https://en.wikipedia.org/wiki/Kesawai_language}
\end{itemize}

\section{Religious and Community Resources}
Joshua Project has a profile for the Kesawai/Namuya people: \url{https://joshuaproject.net/index.php/people_groups/12641}. No scripture or community organisation was found.

\section{Digital and Audio Resources}
No audio, video, corpus or software resources were found.

\section{Curriculum Development Context}
The language has a full modern grammar but no known pedagogical materials. Younger speakers are shifting to Tok Pisin. A curriculum would need new materials drawn from the grammar and community work, for example a practical orthography, story collections and vernacular preschool materials. PNG's elementary vernacular education policy could be a route for this.
\end{document}